We present a causal attention-knockout framework for measuring per-layer
information flow in Vision-Language Models, and apply it to a controlled before/after
pair, Gemma-3-4B and its medical fine-tune MedGemma-4B, on VQA-RAD under a
pre-specified protocol that screens a metric family on a discovery split and re-tests
each candidate on a held-out confirmation split with identical indices for both models.
The pre-specified primary test confirms: relative to its language prior, the medical
model relies less on the question text (question-dominance ratio $5.36 \to 4.15$,
$p{=}0.006$). Four of five discovery candidates replicate: a markedly more decisive
first answer token, smaller closed yes/no logit margins, lower-probability answer text,
and lower image$\to$answer flow in both pathways (the image one at the raw level only). The
medical model is substantially more accurate on closed questions, though the base
model's near-baseline score means the raw gap overstates the true margin. A half-layer
``image enters later'' shift seen
only on discovery does not replicate, which the pre-specified protocol lets us report plainly as
a null rather than a finding.

We deliberately stop short of a causal story. The medical vision encoder differs
sharply from the base (linear CKA 0.155), so the measured decoder routing changes
cannot be attributed to the encoder or the decoder from correlational traces alone.
We frame that attribution as the central open problem and pursue it in follow-up work
via a planned follow-up experiment. The framework is architecture-agnostic and
requires no retraining. The implementation, the analysis scripts, and the per-layer
information-flow curves (as numeric exports, without VQA-RAD text) are released at
\url{https://github.com/HoussemDs/vlm-attention-knockout} to support reproducibility.
